\part{Harmonic Analysis}

\section{Fourier Analysis}

\begin{exercise}[Hilbert spaces and unitary translations]
Work over $\C$.  Define normed and Banach spaces, bounded linear maps and
their operator norms, inner-product spaces, Cauchy sequences, and completion.
Define a Hilbert space to be a complete inner-product space and prove the
projection theorem for a closed subspace.  Deduce the Riesz representation
theorem for bounded linear functionals.  For a bounded operator $T$ define
its adjoint intrinsically by
$\inner{Tx}{y}=\inner{x}{T^*y}$ and prove
$\norm{T^*T}=\norm T^2$.  Define a unitary operator as a surjective linear
isometry and a unitary action as a group homomorphism into such operators.
Construct $L^2(X,\mu)$ from square-integrable functions modulo equality
almost everywhere.  If a locally compact group $G$ acts continuously on $X$
preserving $\mu$, prove that pullback defines a unitary action on
$L^2(X,\mu)$.  Specialize to the additive action of $\R^n$ on itself with
Lebesgue measure and prove strong continuity of its translations.
\end{exercise}

\begin{exercise}[Locally compact groups and Haar measure]
Let $G$ be a locally compact group.  Define a positive linear functional on
$C_c(G)$ and prove its Riesz--Markov representation by outer approximation.
Construct a
nonzero positive left-translation-invariant functional by the covering-number
argument, apply that representation, and prove the uniqueness needed to
obtain left Haar measure $\dd x$, unique up to scale.  Construct right Haar measure and
the continuous homomorphism $\Delta_G:G\to\R_{>0}$ characterized by
\[
 \int_G f(xg)\dd x=\Delta_G(g)^{-1}\int_Gf(x)\dd x.
\]
Define convolution and involution by
\[
 (f*g)(x)=\int_G f(y)g(y^{-1}x)\dd y,
 \qquad f^*(x)=\Delta_G(x^{-1})\overline{f(x^{-1})},
\]
and prove that $L^1(G)$ is a Banach $*$-algebra and that left translation is
unitary on $L^2(G)$.  Prove that $\Delta_G=1$ for abelian and compact groups;
in the abelian case adopt additive notation, so that the formulas reduce to
additive convolution and reflection involution.
\end{exercise}

\begin{exercise}[Pontryagin duality and Fourier decomposition]
Define $\widehat G=\Hom_{\mathrm{cont}}(G,U(1))$ with the compact-open
topology.  Normalize Haar measure on $\widehat G$ so that
\[
 \widehat f(\chi)=\int_G f(x)\overline{\chi(x)}\dd x,
 \qquad
 f(x)=\int_{\widehat G}\widehat f(\chi)\chi(x)\dd\chi
\]
for the functions obtained by convolving two compactly supported continuous
functions.  Extend the transform to a unitary
$L^2(G)\to L^2(\widehat G)$, prove Plancherel and convolution--multiplication,
and prove that evaluation identifies $G$ with $\widehat{\widehat G}$.
Specialize all normalizations to $\R^n$, $\Z^n$, and $\T^n$.
\end{exercise}

\begin{exercise}[Poisson summation and Gaussian Fourier kernels]
On $\R^n$ use
\[
 \widehat f(\xi)=\int_{\R^n}f(x)e^{-2\pi\ii x\cdot\xi}\dd x.
\]
Define the Schwartz space by its seminorms, prove Fourier inversion there,
and compute $\widehat{e^{-\pi t\abs{x}^2}}=t^{-n/2}
e^{-\pi\abs{\xi}^2/t}$.  For a lattice $L$ define its covolume and dual
lattice intrinsically and prove
\[
 \sum_{x\in L}f(x)=\frac1{\operatorname{covol}(L)}
 \sum_{\xi\in L^*}\widehat f(\xi).
\]
Derive the Jacobi transformation
$\vartheta(-1/\tau)=(-\ii\tau)^{1/2}\vartheta(\tau)$ for
$\vartheta(\tau)=\sum_{m\in\Z}e^{\pi\ii m^2\tau}$, fixing the square-root
branch by positivity on the imaginary axis.
\end{exercise}

\begin{exercise}[Dirichlet, Fej\'er, and sinc Fourier kernels]
On $\T$ define
\[
 D_N(\theta)=\sum_{n=-N}^Ne^{\ii n\theta}
 =\frac{\sin((N+1/2)\theta)}{\sin(\theta/2)},\qquad
 F_N=\frac1{N+1}\sum_{j=0}^ND_j .
\]
Prove that $F_N\ge0$, has integral one, and is an approximate identity.
Deduce Fej\'er convergence for continuous functions and the $L^2$ convergence
of Fourier sums.  Define $\operatorname{sinc}u=\sin(u)/u$ continuously at
zero and compute it as the normalized Fourier transform of an interval.
\end{exercise}

\begin{exercise}[Fraunhofer diffraction by rectangular aperture arrays]
For monochromatic scalar light of wavelength $\lambda$, derive from the
Huygens integral in the far-field limit that the amplitude at angle
$\boldsymbol\theta$ is the Fourier transform of the aperture evaluated at
$\boldsymbol\theta/\lambda$.  Deduce the rectangular-aperture intensity as a
product of squared sinc functions.  For $N$ equally spaced identical slits,
derive the squared Dirichlet-kernel array factor, its principal-peak spacing,
and the finite-$N$ peak width.  Identify aperture dimensions, slit spacing,
wavelength, screen distance, and measured intensity as the calibrated and
observed quantities, and state the scalar, paraxial, and far-field
approximations.
\end{exercise}

\section{Peter--Weyl Theory}

\begin{exercise}[Finite-dimensional unitary representations of compact groups]
Define the tensor product $V\otimes W$ by its universal bilinear map and
construct its associativity, unit, and symmetry isomorphisms, writing the
last as $\sigma_{V,W}:V\otimes W\to W\otimes V$.  Call $V$ finite-dimensional
when there are a vector space $V^\vee$ and evaluation
and coevaluation maps
\[
 \operatorname{ev}_V:V^\vee\otimes V\to\C,
 \qquad \operatorname{coev}_V:\C\to V\otimes V^\vee
\]
satisfying the two snake identities.  Define the trace by
\[
\Tr(T)=\operatorname{ev}_V\circ\sigma_{V,V^\vee}\circ
 (T\otimes1)\circ\operatorname{coev}_V
\]
and put $\dim V=\Tr(1_V)$.  Prove cyclicity of trace and its direct-sum and
tensor-product formulas.
Normalize Haar measure on a compact group $K$ to total mass one.  Let
$K\times K$ act on $K$ by
$(a,b)\mathbin\cdot g=agb^{-1}$ and prove that normalized Haar measure is
invariant.  Average an
inner product to prove that every finite-dimensional continuous
representation is unitarizable and every invariant subspace has an invariant
orthogonal complement.  Define irreducibility, intertwiners, contragredients,
and characters.  Prove Schur's lemma and character orthogonality without
choosing coordinates.
\end{exercise}

\begin{exercise}[Peter--Weyl decomposition]
For an irreducible unitary representation $\pi$ define its matrix-coefficient
map
\[
 C_\pi:V_\pi\otimes V_\pi^\vee\longrightarrow C(K),
 \qquad C_\pi(v\otimes\phi)(g)=\phi(\pi(g^{-1})v).
\]
Prove Schur orthogonality and that $\sqrt{\dim V_\pi}\,C_\pi$ is an isometry.
Use convolution, finite-rank convolution projections, and Stone--Weierstrass to
prove the unitary decomposition
\[
 L^2(K)\simeq\widehat\bigoplus_{\pi\in\widehat K}
 V_\pi\otimes V_\pi^\vee,
\]
with left translation on $V_\pi$ and right translation on $V_\pi^\vee$.
Prove the required Stone--Weierstrass density statement from separation of
points and closure under conjugation.  Derive Fourier inversion and
Plancherel for $K$.
\end{exercise}

\begin{exercise}[$SU(2)$, Wigner functions, and characters]
Let $U$ be a two-dimensional Hermitian complex vector space with a unit-norm
nonzero alternating form $\varepsilon:U\otimes U\to\C$.  Define $SU(2)$ as
the unitary automorphisms preserving $\varepsilon$.  On the real Euclidean
space of traceless Hermitian endomorphisms of $U$, with
$\langle A,B\rangle=\tfrac12\Tr(AB)$, prove that conjugation maps
$SU(2)$ onto the connected orthogonal group, denoted $SO(3)$, with kernel
$\{\pm1\}$.  Construct the
permutation coinvariant quotient $\operatorname{Sym}^n(V)$ of $V^{\otimes n}$.
Construct the irreducibles $V_j=\operatorname{Sym}^{2j}(\C^2)$ for
$j\in\frac12\Z_{\ge0}$ and define the Wigner functions
$D^j_{mn}(g)$ as their distinguished weight matrix coefficients.  Derive
their Euler-angle formula, orthogonality, completeness, and
\[
 \chi_j(\theta)=\frac{\sin((2j+1)\theta/2)}{\sin(\theta/2)}
                =U_{2j}(\cos(\theta/2)),
\]
thereby defining the Chebyshev polynomials $U_n$.
\end{exercise}

\begin{exercise}[Clebsch--Gordan and Wigner special functions]
Prove
$V_{j_1}\otimes V_{j_2}\simeq
\bigoplus_{j=\abs{j_1-j_2}}^{j_1+j_2}V_j$.
Define Clebsch--Gordan coefficients as the normalized components of these
canonical isometric inclusions relative to the distinguished weight lines,
and define the Wigner $3j$ and $6j$ symbols from them.  Derive their
orthogonality, the triangle rules, the Racah recoupling identity, and the
integral of three Wigner functions.  Prove the Wigner--Eckart theorem by
identifying tensor operators with intertwiners.
\end{exercise}

\begin{exercise}[Rotational spectroscopy]
Let a rigid symmetric-top molecule have moment-of-inertia operator with
principal values $I_\perp,I_\perp,I_\parallel$.  Quantize its rotational
Hamiltonian by the $SU(2)$ action and prove
\[
 E_{JK}=\frac{\hbar^2}{2I_\perp}J(J+1)
 +\frac{\hbar^2}{2}\left(I_\parallel^{-1}-I_\perp^{-1}\right)K^2.
\]
Express electric-dipole transition amplitudes through Wigner $3j$ symbols and
derive $\Delta J=0,\pm1$, the prohibition $J=0\to J=0$, and the appropriate
$\Delta K$ rule for dipoles parallel or perpendicular to the symmetry axis.
For a linear rotor deduce the line positions
$\nu_{J\to J+1}=\hbar(J+1)/(2\pi I)$ and the relative line strengths from the
$3j$ symbols; these are the quantitative spectroscopic predictions.
\end{exercise}

\section{Symmetric Spaces}

\begin{exercise}[Riemannian symmetric spaces and Gelfand pairs]
Define smooth manifolds by compatible smooth atlases, smooth maps by their
chart representatives, and tangent vectors intrinsically as derivations of
germs.  Construct tangent bundles and define a Lie group as a group object in
smooth manifolds.  Define a Riemannian metric, its volume measure, gradient,
and divergence, and use
$\Delta=-\operatorname{div}\operatorname{grad}$ for the nonnegative
Laplace--Beltrami operator.  Call a connected Riemannian manifold symmetric
when every point $x$ admits an isometry fixing $x$ and acting as $-1$ on
$T_xX$.  Prove that a connected compact symmetric space is a homogeneous
space $G/K$ for the identity component of its isometry group and a compact
stabilizer $K$.

For a compact group $G$ and closed subgroup $K$, define the convolution
algebra of continuous $K$-bi-invariant functions.  Call $(G,K)$ a Gelfand
pair when this algebra is commutative, and prove that this is equivalent to
multiplicity-free decomposition of $L^2(G/K)$.  Define the orthogonal group
$O(V)$ of a Euclidean space $V$.  For $\dim V=d+1$ with $d\geq2$, identify
its unit sphere with $O(V)/O(V)_x$, prove that its geodesic symmetries are induced by the
orthogonal maps fixing $x$ and negating $x^\perp$, and prove that
$(O(d+1),O(d))$ is a Gelfand pair.
\end{exercise}

\begin{exercise}[Gamma and generalized hypergeometric functions]
For $\Re s>0$ define
\[
 \Gamma(s)=\int_0^\infty t^{s-1}e^{-t}\dd t,
 \qquad
 B(a,b)=\int_0^1t^{a-1}(1-t)^{b-1}\dd t,
\]
derive the functional equation and meromorphic continuation of $\Gamma$, and
prove $B(a,b)=\Gamma(a)\Gamma(b)/\Gamma(a+b)$ together with the Euler
reflection and Legendre duplication formulas.  Put
$(a)_0=1$ and $(a)_n=a(a+1)\cdots(a+n-1)$.  When no lower parameter is a
nonpositive integer, define
\[
 {}_pF_q\!\left(\begin{matrix}a_1,\ldots,a_p\\b_1,\ldots,b_q\end{matrix};z\right)
 =\sum_{n\geq0}\frac{(a_1)_n\cdots(a_p)_n}
 {(b_1)_n\cdots(b_q)_n}\frac{z^n}{n!}.
\]
Determine its disk of convergence for $p\leq q$, $p=q+1$, and $p>q+1$, with
terminating exceptions, and derive its differential equation in terms of
$\vartheta=z\partial_z$:
\[
 \left[\vartheta\prod_{j=1}^q(\vartheta+b_j-1)
 -z\prod_{i=1}^p(\vartheta+a_i)\right]y=0.
\]
Derive the confluence limits ${}_2F_1\to{}_1F_1\to{}_0F_1$ by sending an
upper parameter to infinity while rescaling $z$.  Specialize to the Gauss
function ${}_2F_1(a,b;c;z)$ and
derive its three regular singular points, Euler integral, Gauss value at
$z=1$, Euler and Pfaff transformations, and connection formula between the
Frobenius solutions at $0$ and $1$, first when the exponent differences are
nonintegral and then by limits.
\end{exercise}

\begin{exercise}[Jacobi and Gegenbauer functions]
For $\alpha,\beta>-1$ define the Jacobi polynomials by
\[
 P_n^{(\alpha,\beta)}(1-2z)
 =\frac{(\alpha+1)_n}{n!}
 {}_2F_1(-n,n+\alpha+\beta+1;\alpha+1;z).
\]
Derive their differential equation, Rodrigues formula, three-term recurrence,
orthogonality for $(1-t)^\alpha(1+t)^\beta\dd t$, squared norms, and
completeness.  For $\lambda>0$ define the Gegenbauer functions by
\[
 C_n^\lambda(t)=
 \frac{(2\lambda)_n}{(\lambda+1/2)_n}
 P_n^{(\lambda-1/2,\lambda-1/2)}(t)
\]
and derive their generating function.  Recover the Legendre and Chebyshev
families as specializations and
express each $SU(2)$ Wigner small-$d$ function as an elementary half-angle
factor times a Jacobi polynomial, including its normalization.
\end{exercise}

\begin{exercise}[Spherical harmonic decomposition]
Define the degree-$\ell$ polynomial functions on $V$ as
$\operatorname{Sym}^\ell(V^\vee)$ acting by evaluation.  On the unit sphere
$S^d\subset V$, define $\mathcal H_\ell$ as the restrictions of those
annihilated by the Euclidean Laplacian.  Prove intrinsically the orthogonal decomposition
\[
 L^2(S^d)=\widehat\bigoplus_{\ell\geq0}\mathcal H_\ell,
 \qquad
 \Delta_{S^d}|_{\mathcal H_\ell}=\ell(\ell+d-1),
\]
and compute, with a binomial coefficient understood to vanish when its lower
argument is negative,
\[
 N_{d,\ell}=\dim\mathcal H_\ell
 =\binom{\ell+d}{\ell}-\binom{\ell+d-2}{\ell-2}.
\]
Using the Gelfand-pair decomposition, prove that the orthogonal projection
onto $\mathcal H_\ell$ has kernel
\[
 \Pi_\ell(x,y)=\frac{N_{d,\ell}}{\operatorname{vol}(S^d)}
 \frac{C_\ell^{(d-1)/2}(x\cdot y)}
 {C_\ell^{(d-1)/2}(1)}.
\]
Derive the Funk--Hecke formula and its inversion for zonal kernels.  For
$d=2$, define the distinguished spherical harmonics $Y_{\ell m}$ as the
right-$SO(2)$ types among the $SO(3)$ matrix coefficients and recover the
Legendre addition formula without choosing an auxiliary family.
\end{exercise}

\begin{exercise}[Spherical resonator spectroscopy]
Let a scalar wave of propagation speed $v$ be confined to a thin spherical
resonator of radius $R$, with dynamics
$\partial_t^2u+(v^2/R^2)\Delta_{S^d}u=0$.  Use the spherical harmonic
decomposition and Gegenbauer projection kernels to prove that its angular
normal modes have frequencies and degeneracies
\[
 \omega_\ell=\frac vR\sqrt{\ell(\ell+d-1)},
 \qquad \operatorname{mult}(\omega_\ell)=N_{d,\ell}.
\]
For a rotationally symmetric source, express the measured two-point response
as a Gegenbauer series and derive its allowed frequencies and angular nodal
sets.  Separate the scalar thin-shell model, uniform propagation speed, and
perfect rotational symmetry from the measured resonance frequencies,
degeneracies, and angular correlations.
\end{exercise}

\section{Euclidean Group}

\begin{exercise}[Euclidean motion symmetry and momentum orbits]
For a finite-dimensional real inner-product space $V$, define
$E(V)=V\rtimes O(V)$ and its isometric action on the affine Euclidean space
modeled on $V$, and write $d=\dim V$.  Show that Fourier transformation intertwines translations
with multiplication by characters and rotations with their action on $V^*$.
Prove that the nonzero $O(V)$-orbits in $V^*$ are the momentum spheres and
disintegrate Plancherel measure into radial measure and their invariant
probability measures.  Apply the compact-group decomposition of each sphere
to obtain the orthogonal angular-momentum decomposition of $L^2(V)$.
For a transitive locally compact group action on $G/H$ with invariant measure
and a unitary $H$-representation, define the induced representation as the
square-integrable equivariant sections of the associated Hilbert bundle and
prove that left translation is unitary.
Construct the irreducible positive-momentum representations of $E(V)$ by
inducing the corresponding translation character tensored with an
irreducible representation of its orthogonal stabilizer, and identify the
scalar sector with its $O(V)$-fixed induced functions.
\end{exercise}

\begin{exercise}[Bessel functions and Hankel decomposition]
For $c\notin\Z_{\leq0}$ specialize the generalized hypergeometric series to
\[
 {}_0F_1(;c;z)=\sum_{k\geq0}\frac{z^k}{(c)_k k!},
 \qquad
 J_\nu(z)=\frac{(z/2)^\nu}{\Gamma(\nu+1)}
 {}_0F_1(;\nu+1;-z^2/4).
\]
Verify directly the double-confluence limit from the Gauss series.
Derive Bessel's differential equation, its recurrence and Wronskian
relations, and the angular integral that realizes $J_{d/2-1}$ as the
spherical function of the Gelfand pair $(E(d),O(d))$.  Continue $J_\nu$ in
its order and argument, and for $\nu\notin\Z$ define
\[
 Y_\nu(z)=\frac{\cos(\pi\nu)J_\nu(z)-J_{-\nu}(z)}{\sin(\pi\nu)},
 \qquad H_\nu^{(1)}=J_\nu+\ii Y_\nu,
\]
with integral-order values obtained by continuation.  Derive their Wronskian
and incoming/outgoing asymptotics.

Use polar disintegration and the spherical harmonic projections to derive
the radial Fourier formula
\[
 \widehat f(\rho)=2\pi\rho^{1-d/2}
 \int_0^\infty f(r)J_{d/2-1}(2\pi r\rho)r^{d/2}\dd r
\]
for radial Schwartz functions, and the order
$\ell+d/2-1$ Hankel transform on the $\ell$th spherical summand.  Prove its
inversion, Plancherel formula, and completeness kernel from Fourier inversion.
\end{exercise}

\begin{exercise}[Circular-aperture diffraction]
For a circular aperture of diameter $D$ illuminated by monochromatic scalar
light of wavelength $\lambda$, use Euclidean rotational symmetry and the
order-zero Hankel transform to derive
\[
 \frac{I(\theta)}{I(0)}=
 \left[\frac{2J_1(\pi D\sin\theta/\lambda)}
 {\pi D\sin\theta/\lambda}\right]^2.
\]
Prove simplicity of the first positive zero of $J_1$, bound it by convergent
series or Sturm comparison to obtain $j_{1,1}=3.831705\ldots$, and hence
predict the first dark ring
$\sin\theta=1.21967\ldots\lambda/D$.  Derive the encircled-energy profile
from the Bessel recurrence.  Distinguish the predicted angular intensity and
dark-ring radius from aperture calibration, finite screen distance,
polarization, and aberration corrections.
\end{exercise}

\section{Reductive Groups}

\begin{exercise}[Lie algebras, enveloping algebras, and Hopf algebras]
Define a complex Lie algebra, Lie ideals, representations, and the adjoint
representation.  Construct the tensor algebra by its universal property and
define
\[
 U(\mathfrak g)=T(\mathfrak g)/
 \langle x\otimes y-y\otimes x-[x,y]\rangle
\]
by its universal property among associative algebras receiving a Lie map from
$\mathfrak g$.  Filter $U(\mathfrak g)$ by tensor degree, construct the
canonical graded map
$\operatorname{Sym}(\mathfrak g)\to\operatorname{gr}U(\mathfrak g)$, and
prove it is an isomorphism by the rewriting filtration and the Jacobi
identity.  Deduce the Poincar\'e--Birkhoff--Witt theorem, induction and
restriction of modules, and the construction of a Verma module from a
triangular decomposition.  Define coalgebras, bialgebras, and Hopf algebras by
their universal diagrams, and construct the cocommutative Hopf structure on
$U(\mathfrak g)$ for which every $x\in\mathfrak g$ is primitive.
\end{exercise}

\begin{exercise}[Root data and Weyl groups]
Define an affine complex algebraic group by its finitely generated reduced
commutative coordinate Hopf algebra, after defining nilpotent elements and
reduced algebras.  Define connectedness, unipotent groups,
the unipotent radical, and reductivity, and construct its Lie algebra from
derivations at the counit.  For a connected complex reductive group $G$,
define a maximal torus $T$, its
character and cocharacter lattices, roots, coroots, and the Weyl group
$W=N_G(T)/T$.  Prove that $W$ is finite, construct a $W$-invariant Euclidean
metric on the real weight space $X^*(T)\otimes_\Z\R$ by averaging, and prove
that every root reflection is the orthogonal reflection in its root
hyperplane.  Derive the root-space decomposition of $\mathfrak g$ and the
rank-one $\mathfrak{sl}_2$ subalgebra associated with each root.  Choose a
positive system, define simple roots, dominant weights, the Weyl length, and
$\rho$, and prove the chamber and Bruhat decompositions.
\end{exercise}

\begin{exercise}[Highest-weight decomposition]
Define weight spaces without choosing vectors in them.  Prove that every
finite-dimensional irreducible $G$-representation has a unique dominant
highest weight.  For a Borel algebra
$\mathfrak b=\mathfrak t\oplus\mathfrak n_+$ define
$M_\lambda=U(\mathfrak g)\otimes_{U(\mathfrak b)}\C_\lambda$ and construct
the irreducible representation as its unique irreducible quotient.  Use the
rank-one strings to prove finite
dimensionality exactly for dominant integral weights, complete reducibility,
and the highest-weight classification.
\end{exercise}

\begin{exercise}[Weyl characters and Schur functions]
Define exterior powers as alternating tensor quotients and define the
determinant of an endomorphism as the scalar by which it acts on the highest
nonzero exterior power.  Define a partition as a finite weakly decreasing
sequence of nonnegative integers.  In the group ring of the weight lattice
define
$A_\mu=\sum_{w\in W}(-1)^{\ell(w)}e^{w\mu}$.  Construct the BGG resolution
from Verma modules and use its Euler character to prove
\[
 \operatorname{ch}V_\lambda=\frac{A_{\lambda+\rho}}{A_\rho},
 \qquad
 A_\rho=e^\rho\prod_{\alpha>0}(1-e^{-\alpha}).
\]
For $G=GL_r$ identify these characters with the Schur functions
\[
 s_\lambda(x)=
\frac{\det(x_i^{\lambda_j+r-j})}{\det(x_i^{r-j})},
\]
and derive the tableaux expansion, Cauchy identity, and
Littlewood--Richardson tensor-product rule.
\end{exercise}

\begin{exercise}[Flavor multiplets and the Gell-Mann--Okubo prediction]
Use the $SU(3)$ Weyl character $s_{(2,1)}$ to recover the eight weights of the
adjoint representation with their multiplicities.  Let isospin $SU(2)$ and
hypercharge $U(1)$ be the evident subgroup and show that the octet decomposes
as two isodoublets, one isotriplet, and one isosinglet.  Assume the leading
symmetry-breaking mass operator is an isoscalar in $1\oplus8$.  Apply the
Wigner--Eckart theorem and eliminate its three reduced parameters to derive
\[
 2(M_N+M_\Xi)=3M_\Lambda+M_\Sigma.
\]
Explain which four observed baryon centroid masses test the relation and
separate the exact representation-theoretic multiplicities from the
first-order symmetry-breaking assumption.
\end{exercise}

\section{Coadjoint Orbits}

\begin{exercise}[Hamiltonian group actions and moment maps]
For a smooth manifold, construct its cotangent bundle, differential forms,
pullback, wedge product, and exterior derivative, and prove their independence
of charts.  Identify the Lie algebra of a Lie group with its left-invariant
vector fields.  Define symplectic manifolds and
Hamiltonian vector fields by
$\iota_{X_f}\omega=\dd f$, and the Poisson bracket
$\{f,g\}=\omega(X_f,X_g)$.  For a Lie-group action define an equivariant
moment map $\mu:M\to\mathfrak g^*$ by
$\dd\langle\mu,\xi\rangle=\iota_{\xi_M}\omega$.  Prove Noether's theorem,
the Poisson relation for moment-map components, and functoriality under
products.
\end{exercise}

\begin{exercise}[Kirillov--Kostant--Souriau orbits]
For $\lambda\in\mathfrak g^*$ identify
$T_\lambda(G\lambda)$ with $\mathfrak g/\mathfrak g_\lambda$ and define
\[
 \omega_\lambda(\ad^*_\xi\lambda,\ad^*_\eta\lambda)
 =\langle\lambda,[\xi,\eta]\rangle.
\]
Prove that this is well defined, closed, nondegenerate, and has moment map the
orbit inclusion.  Identify the $SU(2)$ orbits with spheres, compute their
symplectic areas intrinsically, and derive the Lie--Poisson equations.
\end{exercise}

\begin{exercise}[Symplectic reduction and orbit quantization]
For a regular value $a$ of a moment map, construct
$M//_aG=\mu^{-1}(a)/G_a$ and prove that its two-form is characterized by
pullback.  Define a prequantum Hermitian line bundle by a unitary connection
of curvature $-\ii\omega/\hbar$.  Prove the integrality condition and, for an
$SU(2)$ orbit of area $2\pi\hbar n$, realize its quantization as the
right-weight-$n$ functions on $SU(2)$ annihilated by the positive-root
operator.  Use Peter--Weyl to identify this space with
$\operatorname{Sym}^n(\C^2)^\vee$.
Derive quantization commutes with reduction for tensor products of these
rank-one orbits and recover the triangle rule.
\end{exercise}

\begin{exercise}[Spin coherent states and their kernels]
Let $J_1,J_2,J_3$ be the infinitesimal $SU(2)$ generators normalized by
$[J_a,J_b]=\ii\hbar\epsilon_{abc}J_c$.  For $j=n/2$, define the normalized
spin coherent state $\lvert x\rangle$ as
the image in $\operatorname{Sym}^n(\C^2)$ of a unit vector projecting to
$x\in S^2$.  Prove
\[
 \abs{\langle x\mid y\rangle}^2=\left(\frac{1+x\cdot y}{2}\right)^{2j},
 \qquad
 \frac{2j+1}{4\pi}\int_{S^2}\lvert x\rangle\langle x\rvert\dd\Omega=1.
\]
Expand a rotated coherent state into the $J_3$ weight lines and show that its
probabilities are the binomial kernel
\[
 p_m(\theta)=\binom{2j}{j+m}
 \cos^{2(j+m)}\frac\theta2\sin^{2(j-m)}\frac\theta2.
\]
Define the Krawtchouk polynomials by
\[
 \sum_{n=0}^N\binom NnK_n(x;p,N)t^n
 =\left(1-\frac{1-p}{p}t\right)^x(1+t)^{N-x}.
\]
Identify the associated rotation coefficients with these polynomials at
$p=\sin^2(\theta/2)$, identify
$K_n(x;p,N)={}_2F_1(-n,-x;-N;1/p)$, and derive their difference equation and
binomial orthogonality.
\end{exercise}

\begin{exercise}[Spin precession and Stern--Gerlach intensities]
For these infinitesimal generators let the Zeeman Hamiltonian be
$H=-\gamma\mathbf B\cdot\mathbf J$.  Solve the Hamiltonian orbit equation and
the quantum evolution to obtain precession frequency
$\omega_L=\abs\gamma\abs{\mathbf B}$.  If a beam polarized along $x$ is
measured along an axis making angle $\theta$ with $x$, prove that the
$2j+1$ separated intensities are precisely the displayed $p_m(\theta)$.  Thus the
coherent-state/Krawtchouk kernel predicts both the line count and every
relative Stern--Gerlach intensity.
\end{exercise}
